\section*{Bab \ref{ch:b3:planetary-interiors-surfaces} --- Interior dan Permukaan Planet}

\begin{solution}{exo:b3:planetary-interiors-surfaces:1}
$\tfrac52C\bar\rho = 2.5\times0.35\times5429 = 4750$; $x^2 = (4750 -
3300)/(5429 - 3300) = 0.681$, $x = 0.825$, inti berjari-jari $0.825\times2440 =
\qty{2010}{km}$; $\rho_c = 3300 + 2129/0.825^3 = \qty{7090}{kg.m^{-3}}$.
\end{solution}

\begin{solution}{exo:b3:planetary-interiors-surfaces:2}
$m = 3500\times\tfrac43\pi(25)^3 = \qty{2.29e8}{kg}$; $E = \tfrac12 m v^2 =
\tfrac12\times\num{2.29e8}\times(\num{2.0e4})^2 = \qty{4.6e16}{J}$, atau
$\num{4.6e16}/\num{4.2e15} = \qty{11}{Mt}$. Aturan ini memberikan kawah berukuran
$20\times\qty{50}{m} = \qty{1}{km}$.
\end{solution}

\begin{solution}{exo:b3:planetary-interiors-surfaces:3}
$N_1 = 120/\num{100000} = \qty{1.2e-3}{km^{-2}}$, galat relatif $1/\sqrt{120}
= \qty{9}{\percent}$. Pada bagian linier, $T = N_1/c = \num{1.2e-3}/\num{4.875e-4}
= \qty{2.46}{Ga}$ (suku eksponensial, $\num{8e-7}$, dapat diabaikan), dan
$\sigma_T = T/\sqrt{120} = \qty{0.22}{Ga}$.
\end{solution}

\begin{solution}{exo:b3:planetary-interiors-surfaces:4}
Luas $4\pi(\num{1.8215e6})^2 = \qty{4.17e13}{m^2}$; daya $2.5\times
\num{4.17e13} = \qty{1.04e14}{W} = \qty{104}{TW}$, 2.2 kali total Bumi
dari permukaan dua belas kali lebih kecil, dan $2.5/0.092 = 27$ kali
fluksnya.
\end{solution}

\begin{solution}{exo:b3:planetary-interiors-surfaces:5}
(a)~$\bar\rho = \num{6.417e23}/(\tfrac43\pi(\num{3.3895e6})^3) =
\qty{3934}{kg.m^{-3}}$; $g = GM/R^2 = \qty{3.73}{m.s^{-2}}$; $\rho gR =
\qty{4.97e10}{Pa}$. (b)~$19\mu/(2\rho gR) = 19\times\num{6.7e10}/\num{9.94e10}
= 12.8$, $k_2 = 1.5/13.8 = 0.109$. (c)~Nilai terukur lebih besar sebesar
2.6 simpangan baku: Mars menghasilkan nilai lebih besar daripada planet padat dengan
kekakuan mantel, seperti yang akan terjadi jika ada inti cair di dalamnya.
\end{solution}

\begin{solution}{exo:b3:planetary-interiors-surfaces:6}
(a)~Bulan: $4\times9.8/1.62 = \qty{24}{km}$ dalam batuan kristalin,
\qty{12}{km} dalam sedimen; Mars dan Merkurius: $4\times9.8/3.7 =
\qty{10.6}{km}$ dan \qty{5.3}{km}. Nilai \qty{20}{km} pada Bulan berada di antara
kedua prediksinya: aturan ini berlaku hingga beberapa puluh persen. (b)~$Y =
\rho gD/3 = 2600\times9.8\times4000/3 = \qty{34}{MPa}$, kekuatan
batuan retak, bukan kristal utuh.
\end{solution}

\begin{solution}{exo:b3:planetary-interiors-surfaces:7}
Mars: $\num{3.22e-14}(e^{6.93T} - 1) + \num{4.875e-4}\,T = \num{2.0e-3}$;
tebakan linier $4.1$~Ga terlalu tua, karena suku eksponensialnya besar
di sana; metode bagi dua memberi $T = \qty{3.34}{Ga}$ (suku eksponensial
$\num{3.6e-4}$, suku linier $\num{1.63e-3}$). Fungsi Bulan memberi
\qty{2.39}{Ga}. Setiap fungsi dikalibrasi untuk satu benda: laju
tumbukan dan ukuran kawah yang dihasilkan suatu proyektil
berbeda antara Bulan dan Mars.
\end{solution}

\begin{solution}{exo:b3:planetary-interiors-surfaces:8}
(a)~$3.551181/1.769138 = 2.0073$ dan $7.154553/3.551181 = 2.0147$, mendekati
$2$. (b)~$0.565247 - 0.844789 + 0.279542 = 0.000000$ per hari, hingga
$\num{1e-7}$: ketiga gerak rata-rata terkunci, suatu resonansi yang menjaga
posisi relatif bulan-bulan itu pada setiap konjungsi. (c)~Tarikan berulang di
titik yang sama pada orbit, sehingga tarikan itu menumpuk alih-alih saling meniadakan;
tarikan itu memasok eksentrisitas secepat pasang surut meredamnya.
\end{solution}

\begin{solution}{exo:b3:planetary-interiors-surfaces:9}
$n = 2\pi/(1.769138\times86400) = \qty{4.111e-5}{s^{-1}}$, $nR =
\qty{74.9}{m.s^{-1}}$, $(nR)^5/G = \qty{3.53e19}{W}$;
$\tfrac{21}{2}e^2(nR)^5/G = 10.5\times\num{1.6e-5}\times\num{3.53e19} =
\qty{5.9e15}{W}$ per satuan $k_2/Q$. Maka $k_2/Q =
\num{1.04e14}/\num{5.9e15} = 0.018$.
\end{solution}

\begin{solution}{exo:b3:planetary-interiors-surfaces:10}
(a)~$x^2 = (\tfrac52C\bar\rho - \rho_m)/(\bar\rho - \rho_m)$ adalah linier terhadap
$C$, dengan kemiringan $\tfrac52\bar\rho/(\bar\rho - \rho_m)$. (b)~Kemiringannya
$2.5\times3340/52 = 161$: $x^2$ berubah sebesar $0.161$ untuk $0.001$ pada $C$.
Untuk $C = 0.395$, $x^2 = 0.036 + 0.161 = 0.197$, $x = 0.44$, $\rho_c =
\qty{3880}{kg.m^{-3}}$; untuk $C = 0.393$, $x^2 < 0$: tidak ada inti sama sekali.
(c)~Dekat $C = \tfrac25$ momen inersia hampir tak peka terhadap inti kecil:
angka penting ketiga menentukan apakah ada inti besar yang ringan atau tidak sama sekali.
Diperlukan gelombang seismik, pasang surut, dan medan magnet.
\end{solution}

\begin{solution}{exo:b3:planetary-interiors-surfaces:11}
(a)~$\sigma_T = N_1/(10\,N_1')$: pada \qty{1.0}{Ga}, $N_1 \approx N_1' \approx
\num{4.875e-4}$, $\sigma_T = \qty{0.10}{Ga}$; pada \qty{3.8}{Ga}, $N_1 =
\num{1.07e-2}$ dan $N_1' = \num{6.15e-2}$, $\sigma_T = \qty{0.017}{Ga}$,
enam kali lebih kecil. (b)~Dengan $N_1 = cT$, $N_1' = c$ dan $\sigma_T =
cT/(\sqrt n\,c) = T/\sqrt n$. (c)~Fungsi kronologi itu sendiri
tidak pasti: ia bertumpu pada beberapa lokasi di Bulan yang umurnya telah diketahui dan pada
penskalaan dari Bulan ke Mars, galat yang tidak dapat dikurangi dengan menghitung lebih banyak kawah.
\end{solution}

\begin{solution}{exo:b3:planetary-interiors-surfaces:12}
(a)~Rasionya adalah $(n_ER_E/n_IR_I)^5(e_E/e_I)^2 = (0.4982\times0.8569)^5
\times(0.009/0.004)^2 = 0.0142\times5.06 = 0.072$: kira-kira
\qty{7.5e12}{W}, atau $\num{7.5e12}/(4\pi(\num{1.5608e6})^2) =
\qty{0.24}{W.m^{-2}}$. (b)~\omterm{def:b3:planetary-interiors-surfaces:ocean-world}{Samudra bawah permukaan} memisahkan cangkang es dari
bagian dalam: cangkang melentur jauh lebih banyak daripada bulan padat, dan esnya
mendisipasi energi secara berbeda dari batuan panas Io, sehingga $k_2$ dan $Q$ keduanya berubah.
\end{solution}

\begin{solution}{pb:b3:planetary-interiors-surfaces:1}
\textbf{1.} $\bar\rho = \num{6.4169e23}/(\tfrac43\pi(\num{3.3895e6})^3) =
\qty{3934}{kg.m^{-3}}$.

\textbf{2.} $\tfrac52C\bar\rho = 3600$; $x^2 = (3600 - 3500)/(3934 - 3500) =
0.230$, $x = 0.479$, inti berjari-jari \qty{1620}{km}; $\rho_c = 3500 +
434/0.479^3 = \qty{7450}{kg.m^{-3}}$.

\textbf{3.} $x = 1830/3389.5 = 0.540$; $\rho_m = (3600 - 0.2915\times
3934)/(1 - 0.2915) = \qty{3460}{kg.m^{-3}}$, $\rho_c = 3462 +
472/0.1574 = \qty{6460}{kg.m^{-3}}$.

\textbf{4.} Nilainya sepertiga lebih rendah. Tekanan di inti Mars jauh
lebih rendah daripada di inti Bumi, sehingga logamnya kurang termampatkan; dan
inti yang sangat ringan menunjukkan bahwa unsur-unsur jauh lebih ringan bercampur dengan besi.

\textbf{5.} \omterm{def:b3:planetary-interiors-surfaces:love-number}{Bilangan Love} (2003) muncul jauh sebelum InSight (2021): ia
sudah menunjukkan bahwa setidaknya inti luar itu cair dan menempatkan jari-jarinya
antara 1520 dan \qty{1840}{km}, rentang yang memuat
nilai seismologis.

\textbf{6.} $1 - (1 - 20/3389.5)^3 = \qty{1.8}{\percent}$ hingga $1 - (1 -
37/3389.5)^3 = \qty{3.2}{\percent}$ dari volumenya, dan lebih kecil dari massanya:
lebih ringan daripada mantel, ia menurunkan massa jenis mantel efektif sedikit,
pengaruhnya lebih kecil daripada ketidakpastian model.

\textbf{7.} Dari \qty{1}{km}: 54; dari 1.41: 16; dari 2: 10; dari 2.83: 5;
dari 4, 5.66, 8 dan 11.3: 2; dari 16: 1.

\textbf{8.} $N = 54/80000 = \qty{6.75e-4}{km^{-2}}$, $\sigma = \sqrt{54}/80000
= \qty{0.92e-4}{km^{-2}}$ (\qty{14}{\percent}).

\textbf{9.} Diharapkan $54\times0.355 = 19$ kawah dari \qty{1.41}{km} dan
$54\times0.147 = 7.9$ dari \qty{2}{km}; teramati 16 dan 10, dalam satu
simpangan baku ($\sqrt{19} = 4.4$, $\sqrt{7.9} = 2.8$): konsisten.

\textbf{10.} Tidak: $0.072/\num{6.75e-4} = 107$ kali lebih rendah darinya. Hitungan
mencatat waktu.

\textbf{11.} Interval dengan lebar sama pada $\log D$ meliputi rentang itu secara merata dan
menjaga jumlah kawah yang sebanding di setiap interval pada distribusi yang curam.

\textbf{12.} Itu kawah sekunder, dibentuk oleh bongkahan yang terlempar dari
kawah primer berukuran besar: bukan tumbukan independen dari angkasa, sehingga akan menambah
kawah yang tidak dihitung oleh fungsi kronologi.

\textbf{13.} $T = \num{6.75e-4}/\num{4.875e-4} = \qty{1.385}{Ga}$; suku
eksponensial, $\num{3.22e-14}(e^{9.6} - 1) = \num{4.7e-10}$, adalah
sepersejuta dari $N_1$.

\textbf{14.} $\sigma_T = T/\sqrt{54} = \qty{0.19}{Ga}$: $1.38 \pm 0.19$~Ga.

\textbf{15.} $\num{6.75e-4}/\num{8.38e-4} = \qty{0.81}{Ga}$, lebih muda sebesar
hampir setengahnya.

\textbf{16.} $31 + 12 + 7 + 2 + 0 + 2 + 2 = 56$ kawah; $N(\qty{2}{km}) =
56/120000 = \qty{4.67e-4}{km^{-2}}$; $N_1 = \num{4.67e-4}/0.1473 =
\qty{3.17e-3}{km^{-2}}$. Dengan metode bagi dua pada $\num{3.22e-14}(e^{6.93T} - 1) +
\num{4.875e-4}\,T = \num{3.17e-3}$ menghasilkan $T = \qty{3.54}{Ga}$ (suku
eksponensial \num{1.45e-3}, suku linier \num{1.73e-3}).

\textbf{17.} $N_1' = \num{3.22e-14}\times6.93\,e^{6.93\times3.54} +
\num{4.875e-4} = \num{1.05e-2}$; $\sigma_T = \num{3.17e-3}/(\sqrt{56}\times
\num{1.05e-2}) = \qty{0.040}{Ga}$.

\textbf{18.} Dataran A: \qty{14}{\percent}; dataran B: \qty{1.1}{\percent}. Dalam
bagian awal kronologi yang curam, perubahan besar massa jenis kawah
berarti perubahan kecil pada umur; pada bagian linier, umur sama
tidak pastinya dengan hitungannya.

\textbf{19.} $120000\times\num{3.17e-3}\times(1 - 0.1473) = 324$; ini
memuat $48 + 8 = 56$.

\textbf{20.} $N_1 = 56/(120000\times0.8527) = \qty{5.47e-4}{km^{-2}}$,
$T_r = \num{5.47e-4}/\num{4.875e-4} = \qty{1.12}{Ga}$.

\textbf{21.} $\sigma = 1.12/\sqrt{56} = \qty{0.15}{Ga}$.

\textbf{22.} Kawah \qty{2}{km} berkedalaman \qty{400}{m}: lava
setidaknya setebal itu. Kawah yang lebih besar lebih dalam dan tepinya berdiri lebih tinggi;
sebagian telah terisi, tetapi masih tampak.

\textbf{23.} $1.38 \pm 0.19$ dan $1.12 \pm 0.15$~Ga berbeda sebesar \qty{0.26}{Ga},
sekitar satu simpangan baku gabungan ($\sqrt{0.19^2 + 0.15^2} = 0.24$):
hitungannya memungkinkan satu episode tunggal, meskipun belum membuktikannya.

\textbf{24.} Dataran lava muda A memiliki \omterm{def:b3:advanced-geochronology:model-age}{umur model}
\textbf{1.4~miliar~tahun} ($1.38 \pm 0.19$~Ga).
\end{solution}
